
\section{Implementation Blueprint}


\begin{figure}[H]
\centering
\begin{tikzpicture}[node distance=0.42cm, every node/.style={figlabel}, layer/.style={figbox, minimum width=9.6cm, minimum height=.74cm, align=center}]
\path[figpanel] (-0.85,-0.65) rectangle (10.45,4.55);
\node[figtitle, anchor=west] at (-0.4,4.08) {CONTROL STACK};
\node[layer, fill=BitGray] (p) at (4.8,0) {Markdown state};
\node[layer, above=of p, fill=BitBlue!8, draw=BitBlue!65] (s) {Schema};
\node[layer, above=of s, fill=BitTeal!8, draw=BitTeal!70] (l) {Linter};
\node[layer, above=of l, fill=BitGold!10, draw=BitGold!75!black] (h) {Hook / continuous integration};
\node[layer, above=of h, fill=BitRose!10, draw=BitRose!65] (t) {Test / sandbox / approval};
\draw[figarrow] ($(p.east)+(0.35,0)$) -- node[right, figsmall] {harder control} ($(t.east)+(0.35,0)$);
\end{tikzpicture}
\caption{Readable text is useful. Reliable control requires stronger layers around it.}
\label{fig:enforcement-stack}
\end{figure}


A practical Continuity implementation has two layers. The first layer is human-readable Markdown or YAML front matter. The second layer is deterministic validation.

\subsection{Example file header}

\begin{lstlisting}
---
continuity_type: state
scope: repository
loads_on: boot
mutability: overwrite-current
updates_after:
  - architecture_change
  - task_completion
  - failed_assumption
validation:
  - max_lines: 180
  - no_unresolved_placeholders
  - no_conflicting_status
---
\end{lstlisting}

\subsection{Hooks and validators}

A useful implementation should provide at least five checks.
\begin{enumerate}
  \item \textbf{Pre-load budget checks} stop archive files or oversized manifests from loading by default.
  \item \textbf{Pre-action scope checks} verify that the agent has loaded the nearest relevant state chain.
  \item \textbf{Post-action reconciliation checks} flag when code changed but \statefile{STATE.md} or \statefile{TODO.md} did not.
  \item \textbf{Artifact checks} reject task completion without evidence links, test output, or commits.
  \item \textbf{Safety checks} detect nested files that weaken parent safety rules.
\end{enumerate}

\subsection{Promotion rules}

Generated state should not become canonical automatically. A run summary written by an agent can seed \file{LEARNINGS.md} or \file{completed/}, but promotion into \statefile{STATE.md} or repository-wide memory should pass through review, linting, and freshness checks.

\subsection{Why this layer matters}

The implementation layer prevents the central category error of the whole topic: confusing readable text with reliable control. Continuity is strongest when the repository shows the state and the surrounding tooling checks that the state still makes sense.
